\section{An explicit level for the existential fragment}
\label{sec:existential-bound}

For an ordinary existential sentence, the polynomial to be tested is
explicitly listed after the quartic reduction.  This permits more precise
level accounting than the controlled-family formulation requires.
There is also a simpler root selection: we need to isolate some root below
a threshold, whereas the two-block argument needs to specify any prescribed
coordinate root.  The simpler selection yields the following refinement.

\begin{theorem}\label{thm:etr-bound}
For the circuit encoding of the existential theory of the reals used in
this paper,
\[
 \ETR\in\Clevel{26},
 \qquad \existsR\subseteq\Clevel{26}.
\]
\end{theorem}

We prove the bound using the preceding algebraic and counting tools,
with a short list of derivative conditions in place of the labels from
Section~\ref{sec:labels}.

\subsection{Selecting some root with few derivative conditions}

The distinction between selecting some root and selecting a prescribed
root is elementary but useful.  Derivative signs distinguish all real
roots by Lemma~\ref{lem:thom-separation}.  If only one root from a
nonempty set is needed, repeatedly keeping a smallest nonempty sign
class gives a much shorter description.

\begin{lemma}[Sparse root selection]\label{lem:sparse-selection}
Let $p\in\R[y]$ be a nonzero polynomial of degree $d\ge1$, and let
$S$ be a nonempty set of distinct real roots of $p$.  There is a list
of at most $\lceil\log_2|S|\rceil$ conditions
\[
 \sgn p^{(j)}(\alpha)=v,
 \qquad 1\le j\le d,\quad v\in\{-1,0,1\},
\]
satisfied by exactly one member of $S$.
\end{lemma}

\begin{proof}
If at least two roots remain, their derivative sign strings differ in
some coordinate.  That coordinate partitions the remaining roots into
at least two nonempty classes.  Record the condition defining a smallest
nonempty class and keep that class; its size is at most half the previous
size.  After at most $\lceil\log_2|S|\rceil$ repetitions exactly one
root remains.  If $|S|=1$, use the empty list.
\end{proof}

For an $\ETR$ instance $\exists y\,\Theta(y)$, apply
Lemma~\ref{lem:quartic-fibers} with no free parameters and with
$\Phi=\neg\Theta$.  Renaming the full variable tuple gives an explicitly
listed nonnegative polynomial $F\in\Z[x_1,\ldots,x_n]$ with
$n\ge1$ and degree at most four, whose real zeros encode truth of
the instance.
Use Proposition~\ref{prop:specialization} with $\ell=5$.
In its notation, $D=5^n$, $U=2^\sigma$, $R=2^\rho$, and
\[
 q(y)=\det(yI-M_{h_U})\in\Z[y]
\]
is monic and squarefree of degree $D$.  The original sentence is true
if and only if $q$ has a real root below $R$.
Take $K=\max(1,H+\sigma V)$ from that proposition, so
$\norm q\le2^K$ and $K$ is short and polynomial-time computable.
Although $U$, $R$, and the coefficients
of $q$ may have long binary expansions, their bound exponents are short.

Put $T=2+3n$.  A test list consists of $t$ pairs $(b_i,v_i)$, where
\[
 1\le t\le T,\qquad b_1(y)=R-y,\qquad v_1=1,
\]
and each subsequent $b_i$ is $q^{(j_i)}$ for some $1\le j_i\le D$,
with $v_i\in\{-1,0,1\}$.  A root matches the list when
$\sgn b_i(\alpha)=v_i$ for every $i$.  If a real root below $R$
exists, Lemma~\ref{lem:sparse-selection} applied to all such roots
gives a list matching exactly one real root.  Indeed, at most
$\lceil\log_2D\rceil\le3n$ derivative conditions suffice.
Every list has polynomial encoding length, since both $T$ and
$\log D$ are polynomially bounded.

\subsection{A polynomial detecting the list}

We next turn these conditions into a polynomial that is negative exactly
at matching real roots.  The construction is the rational sign
approximation of Lemma~\ref{lem:rational-sign}, now applied to only
polynomially many tests.

First we need a common range for their nonzero values.
Every complex root $\alpha$ of $q$ satisfies $|\alpha|\le2^{K+1}$,
and every test polynomial satisfies
\[
 \deg b_i\le D,
 \qquad \norm{b_i}\le2^{K+D^2+\rho+2}.
\]
For derivatives, the coefficient increase is bounded by
$D^D\le2^{D^2}$; the same displayed estimate covers $R-y$.
Set
\begin{equation}\label{eq:sparse-value-bounds}
 B_0=K+D^2+\rho+2+D(K+1),\qquad B=DB_0+2.
\end{equation}
Then $|b_i(\alpha)|\le2^{B_0}$.
For the lower bound, the eigenvalues of the integer matrix
$b_i(M_{h_U})$ are these values at the $D$ roots of $q$.
The product of its nonzero eigenvalues, up to sign, is a nonzero
coefficient of its monic integer characteristic polynomial, and so
has absolute value at least one.  Bounding the other factors by
$2^{B_0}$ gives
\begin{equation}\label{eq:sparse-value-range}
 b_i(\alpha)\ne0
 \quad\Longrightarrow\quad
 2^{-B}\le |b_i(\alpha)|\le2^B.
\end{equation}
This argument permits some test values to vanish.

Let $s=T+10$ and define
\begin{equation}\label{eq:sparse-approximant}
 P(z)=\prod_{j=-B}^{B}
 \left(2^{\max(0,j)}+2^{\max(0,-j)}z\right)^{2s},
 \qquad
 A(z)=\frac{P(z)-P(-z)}{P(z)+P(-z)}.
\end{equation}
By Lemma~\ref{lem:rational-sign}, the denominator is at least one
on the real line, $A(0)=0$, and Equation~\eqref{eq:sparse-value-range}
implies
\[
 |A(b_i(\alpha))-\sgn b_i(\alpha)|
 \le2\cdot3^{-2s}\le\frac1{8T}
 \qquad(q(\alpha)=0,\ \alpha\in\R).
\]
Consequently
\[
 \Psi(y)=4\sum_{i=1}^{t}(A(b_i(y))-v_i)^2-1
\]
is at most $-1+1/(16T)$ at a matching real root.  At any other
real root one difference in the sum has absolute value at least
$1-1/(8T)\ge7/8$, and therefore $\Psi\ge33/16$.

Clear the positive squared denominators by setting
\begin{align*}
 D_i(y)&=P(b_i(y))+P(-b_i(y)),\\
 N_i(y)&=P(b_i(y))-P(-b_i(y)),\\
 g_*(y)&=4\sum_{i=1}^{t}(N_i(y)-v_iD_i(y))^2
                 \prod_{j\ne i}D_j(y)^2
              -\prod_{j=1}^{t}D_j(y)^2.
\end{align*}
This is an integer polynomial, equal on the real line to
$\Psi\prod_iD_i^2$.  As every $D_i\ge1$, it satisfies
\begin{equation}\label{eq:sparse-sign-margin}
 \begin{cases}
  g_*(\alpha)\le-\tfrac12,&\text{if $\alpha$ matches the list},\\
  g_*(\alpha)\ge\tfrac12,&\text{otherwise},
 \end{cases}
 \qquad(q(\alpha)=0,\ \alpha\in\R).
\end{equation}
For one additional choice $e\in\{0,\ldots,D\}$ put
\begin{equation}\label{eq:sparse-product}
 g(y)=4Dg_*(y)+e,
 \qquad
 X=\det M_{g(h_U)}
   =\prod_{q(\alpha)=0}g(\alpha)\in\Z.
\end{equation}
The determinant identity is Lemma~\ref{lem:real-eigenline}.
The sign argument of Section~\ref{sec:existential}, applied to
$L=-g_*$, gives
\begin{equation}\label{eq:sparse-product-criterion}
 q\text{ has a real root below }R
 \quad\Longleftrightarrow\quad
 X<0\text{ for some test list and some }e.
\end{equation}
Indeed, every $e$ preserves the real-root signs.  A list matching
exactly one real root gives exactly one negative real factor;
some $e$ avoids every nonreal zero, since each of at most $D$
nonreal roots forbids at most one of the $D+1$ choices.
The conjugate pairs then contribute positive factors.
Conversely, a negative product has a negative real factor, whose
root matches the list and hence lies below $R$.

For the subsequent computations we record explicit short bounds.
Writing $m_*=2s(2B+1)$ for the degree of $P$, set
\begin{equation}\label{eq:sparse-size-bounds}
 \begin{split}
 d_g&=2Tm_*D,\\
 l_*&=m_*(2B+1),\\
 E_*&=2T(l_*+2)+16(D+T+2),\\
 m_X&=DE_*+1.
 \end{split}
\end{equation}
Then $\deg g\le d_g$ and $|X|\le2^{m_X}$.
For the first bound, every term in $g_*$ has $t$ squared factors,
each of degree at most $m_*D$.  For the second, at every complex root
$\alpha$ each base in $P(\pm b_i(\alpha))$ has modulus at most
$2^{2B+1}$.  Hence
\[
 |P(\pm b_i(\alpha))|\le2^{l_*},\qquad
 |D_i(\alpha)|,|N_i(\alpha)|\le2^{l_*+1},\qquad
 |N_i(\alpha)-v_iD_i(\alpha)|\le2^{l_*+2}.
\]
Therefore
\[
 |g(\alpha)|\le D(16T+5)\,2^{2T(l_*+2)}\le2^{E_*},
\]
which gives the asserted product bound.  All displayed bound
parameters use a fixed number of sums and products of short integers,
so are themselves short and polynomial-time computable.

\subsection{Accounting for the oracle levels}

The two multiplication-determinant computations now have an explicit
cost.  Set
\[
 \Omega=6(1+d_g),\qquad L_0=D\Omega+4n,\qquad
 \lambda=5L_0+4n+D^2+10.
\]
For every prime $p>\lambda$, Proposition~\ref{prop:char-coeff}
applies both to $h_U$, of degree at most six, and to $g(h_U)$,
of degree at most $6d_g$.  The larger cutoff is used for both
applications.  To recover $\sgn X$ we need only
$p\le Y=2\max(16,m_X,\lambda)^4$, as in
Proposition~\ref{prop:crt-sign}.  In particular every queried prime
has polynomial bit length.

For the first application, the relations $f_i(\cdot,U)$ and
the objective $h_U$ have direct polynomial-bit-time evaluations
modulo $p$.  Indeed, $\ell=5$ gives $a_i=20^i$, with polynomial
bit length.  The quartic $F$ and its first and second derivatives
have polynomial-size monomial lists.  The constants
$U\bmod p$ and $R\bmod p$ are computed by powering with the short
exponents $\sigma$ and $\rho$.  The Jacobian is an $n\times n$
determinant of directly evaluable entries.  Thus
Proposition~\ref{prop:char-coeff}, with $a=0$, computes the indexed
coefficients $q_k\bmod p$ of $q(y)=\sum_{k=0}^{D}q_ky^k$ in
$\Flevel8$.

For $1\le j\le D$,
\[
 q^{(j)}(y)=\sum_{k=j}^{D}q_k
       \left(\prod_{r=0}^{j-1}(k-r)\right)y^{k-j}.
\]
The falling factorial is in $\Flevel4$ by Lemma~\ref{lem:product};
the outer sum therefore gives derivative evaluations in
$\Flevel9$ by Lemma~\ref{lem:counting}.  This includes every
test $b_i$, since $R-y$ is directly evaluable.
Each factor of $P(b_i(y))$ in
Equation~\eqref{eq:sparse-approximant} is consequently in
$\Flevel9$.  The product lemma puts $P(\pm b_i(y))$ in
$\Flevel{13}$.  Since the list contains only polynomially many
tests, the formulas for $D_i,N_i,g_*$, and $g$ need only
polynomially many such calls and short arithmetic operations.
Thus $g(y)$, and then $g(h_U(x))$, are evaluable in
$\Flevel{13}$.  These expressions use only integer additions and
multiplications; their field evaluations never invert a test denominator.

The second application of Proposition~\ref{prop:char-coeff}, with
$a=13$ and $v=g(h_U)$, computes
\[
 X\bmod p=(-1)^D[y^0]\det(yI-M_{g(h_U)})\bmod p
\]
in $\Flevel{21}$.  The first application uses only the original
quartic and its critical equations; the second calls the first
through the evaluations of $g$.  This is a fixed sequence of two
applications, independent of $n$ and $D$.

By Proposition~\ref{prop:crt-sign} and
Equation~\eqref{eq:sparse-size-bounds}, the sign of $X$ is in
$\Flevel{24}$.  Lemma~\ref{lem:counting} then puts existence of
a test list and an $e$ with $X<0$ in $\Flevel{25}$.
Equation~\eqref{eq:sparse-product-criterion} therefore gives a
decision function for $\ETR$ at that level, and a language in
$\Clevel{26}$.  Polynomial-time many-one precomputation preserves
this level, proving the second assertion of Theorem~\ref{thm:etr-bound}.

\begin{table}[tb]
\centering
\begin{tabular}{lc}
\toprule
Quantity or decision & Sufficient class\\
\midrule
Indexed coefficients of $q\bmod p$ & $\Flevel8$\\
Derivative-test evaluations & $\Flevel9$\\
Evaluations of $P$, $g$, and $g(h_U)$ & $\Flevel{13}$\\
$X\bmod p$ & $\Flevel{21}$\\
$\sgn X$ & $\Flevel{24}$\\
Existence of a list and $e$ with $X<0$ & $\Flevel{25}$\\
Membership in $\ETR$ & $\Clevel{26}$\\
\bottomrule
\end{tabular}
\caption{A conservative level bound for the existential fragment.
Every indexed family uses a single uniform evaluator.}
\label{tab:existential-levels}
\end{table}

The shortness of all indices is independent of these level counts.
The bounds $D,K,\rho,\sigma,B,d_g,m_X,\lambda,Y$ are computed
in polynomial bit time, test lists have at most $T$ entries, and
every field element and product index has a polynomial-length
encoding.  Thus no long integer is passed as an explicit oracle
query, and Table~\ref{tab:existential-levels} describes fixed uniform
machines.
